\subsubsection*{Completeness of the Standard Connectives}

% BEGIN GENERATED ARTIFACT: thm:standard-connectives-functionally-complete-propositional-logic
\begin{theorembox}{Theorem (Standard Connectives are Functionally Complete)}
\begin{theorem}[Standard Connectives are Functionally Complete]
\label{thm:standard-connectives-functionally-complete-propositional-logic}
The connective basis
\[
\{\neg,\land,\lor\}
\]
is functionally complete.

\noindent\hyperref[prf:standard-connectives-functionally-complete-propositional-logic]{Go to proof.}
\end{theorem}
\end{theorembox}

\begin{remark*}[Standard quantified statement]
\[
\operatorname{FunctionallyComplete}(\{\neg,\land,\lor\}).
\]
\end{remark*}

\begin{remark*}[Predicate reading]
Every Boolean function is represented by a formula using only negation,
conjunction, and disjunction.
\end{remark*}

\begin{remark*}[Interpretation]
Canonical disjunctive normal form uses only literals, conjunctions, and
disjunctions. Since literals use variables and negation, the normal-form theorem
shows that \(\{\neg,\land,\lor\}\) can express every Boolean function.
\end{remark*}

\begin{dependencies}
\begin{itemize}
  \item \hyperref[def:functionally-complete-connective-set-propositional-logic]{Functionally Complete Connective Set}
  \item \hyperref[thm:dnf-representation-of-boolean-functions-propositional-logic]{DNF Representation of Boolean Functions}
  \item \hyperref[def:disjunctive-normal-form-propositional-logic]{Disjunctive Normal Form}
\end{itemize}
\end{dependencies}
% END GENERATED ARTIFACT: thm:standard-connectives-functionally-complete-propositional-logic

% BEGIN GENERATED ARTIFACT: thm:negation-conjunction-functionally-complete-propositional-logic
\begin{theorembox}{Theorem (Negation and Conjunction are Functionally Complete)}
\begin{theorem}[Negation and Conjunction are Functionally Complete]
\label{thm:negation-conjunction-functionally-complete-propositional-logic}
The connective basis
\[
\{\neg,\land\}
\]
is functionally complete.

\noindent\hyperref[prf:negation-conjunction-functionally-complete-propositional-logic]{Go to proof.}
\end{theorem}
\end{theorembox}

\begin{remark*}[Standard quantified statement]
\[
\operatorname{FunctionallyComplete}(\{\neg,\land\}).
\]
\end{remark*}

\begin{remark*}[Predicate reading]
Every Boolean function is represented by a formula using only negation and
conjunction.
\end{remark*}

\begin{remark*}[Interpretation]
Disjunction is definable from negation and conjunction by De Morgan's law:
\[
\varphi\lor\psi\equiv\neg(\neg\varphi\land\neg\psi).
\]
Thus every formula using \(\neg,\land,\lor\) can be translated into an equivalent
formula using only \(\neg\) and \(\land\).
\end{remark*}

\begin{dependencies}
\begin{itemize}
  \item \hyperref[thm:standard-connectives-functionally-complete-propositional-logic]{Standard Connectives are Functionally Complete}
  \item \hyperref[thm:de-morgan-laws-propositional-logic]{De Morgan Laws for Propositional Logic}
  \item \hyperref[thm:replacement-of-logical-equivalents-propositional-logic]{Replacement of Logical Equivalents}
\end{itemize}
\end{dependencies}
% END GENERATED ARTIFACT: thm:negation-conjunction-functionally-complete-propositional-logic

% BEGIN GENERATED ARTIFACT: thm:negation-disjunction-functionally-complete-propositional-logic
\begin{theorembox}{Theorem (Negation and Disjunction are Functionally Complete)}
\begin{theorem}[Negation and Disjunction are Functionally Complete]
\label{thm:negation-disjunction-functionally-complete-propositional-logic}
The connective basis
\[
\{\neg,\lor\}
\]
is functionally complete.

\noindent\hyperref[prf:negation-disjunction-functionally-complete-propositional-logic]{Go to proof.}
\end{theorem}
\end{theorembox}

\begin{remark*}[Standard quantified statement]
\[
\operatorname{FunctionallyComplete}(\{\neg,\lor\}).
\]
\end{remark*}

\begin{remark*}[Predicate reading]
Every Boolean function is represented by a formula using only negation and
disjunction.
\end{remark*}

\begin{remark*}[Interpretation]
Conjunction is definable from negation and disjunction by De Morgan's law:
\[
\varphi\land\psi\equiv\neg(\neg\varphi\lor\neg\psi).
\]
Thus every formula using \(\neg,\land,\lor\) can be translated into an equivalent
formula using only \(\neg\) and \(\lor\).
\end{remark*}

\begin{dependencies}
\begin{itemize}
  \item \hyperref[thm:standard-connectives-functionally-complete-propositional-logic]{Standard Connectives are Functionally Complete}
  \item \hyperref[thm:de-morgan-laws-propositional-logic]{De Morgan Laws for Propositional Logic}
  \item \hyperref[thm:replacement-of-logical-equivalents-propositional-logic]{Replacement of Logical Equivalents}
\end{itemize}
\end{dependencies}
% END GENERATED ARTIFACT: thm:negation-disjunction-functionally-complete-propositional-logic
